%% Supplementary material for the IEEE TPAMI submission (IEEEtran, compsoc).
%% Build main.tex and this file twice each: they read each other's labels
%% through the xr package.
\documentclass[10pt,journal,compsoc]{IEEEtran}

% Core mathematics and typography
\usepackage{amsmath,amssymb,amsfonts}
\usepackage{amsthm}
\usepackage{mathtools}
\usepackage{graphicx}
% Font expansion and protrusion keep justified lines from opening wide gaps.
\usepackage{microtype}
\usepackage{booktabs}
\usepackage{multirow}
\usepackage{array}
\usepackage{placeins}
% Labels of the main paper, for references such as Eq. (2).
\usepackage{xr}
\externaldocument[M-]{main}
\usepackage{supertabular}
% A table caption set in the text flow, in the style of IEEEtran's floating
% tables, for a breakable supertabular that follows it.
\makeatletter
\newcommand{\flowtablecaption}[1]{\par\begingroup\@parboxrestore\normalsize
  \def\@captype{table}\refstepcounter{table}%
  \@makecaption{\fnum@table}{#1}\par\endgroup
  \vspace{-5pt}}% the baseline skip before the table replaces part of the float gap
\makeatother
% IEEEtran's compsoc headings, with the same natural spacing but less stretch,
% so that a short column spreads its spare space over paragraph breaks and
% displays instead of opening wide gaps above headings.
\makeatletter
\def\section{\@startsection{section}{1}{\z@}{-3.5ex plus -0.5ex minus -1.5ex}%
{0.7ex plus 0.2ex minus 0ex}{\normalfont\sublargesize\sffamily\bfseries\scshape}}
\def\subsection{\@startsection{subsection}{2}{\z@}{-3.5ex plus -0.5ex minus -1.5ex}%
{0.7ex plus 0.1ex minus 0ex}{\normalfont\normalsize\sffamily\bfseries}}
\makeatother
\usepackage{cite}
\usepackage{xcolor}
\usepackage{url}

% Uncomment when algorithms are needed.
% \usepackage{algorithm}
% \usepackage{algorithmic}

% Uncomment for hyperlinks in the review/preprint version.
% \usepackage[hidelinks]{hyperref}

% Common theorem environments
\newtheorem{theorem}{Theorem}
\newtheorem{lemma}[theorem]{Lemma}
\newtheorem{proposition}[theorem]{Proposition}
\theoremstyle{definition}
\newtheorem{definition}{Definition}

% Convenient notation
\newcommand{\R}{\mathbb{R}}
\newcommand{\vect}[1]{\mathbf{#1}}
\DeclareMathOperator*{\argmax}{arg\,max}
\newcolumntype{L}[1]{>{\raggedright\arraybackslash}p{#1}}

% Number sections, tables and equations with an S prefix.
\renewcommand{\thesection}{S\arabic{section}}
\renewcommand{\thetable}{S\arabic{table}}
\renewcommand{\thefigure}{S\arabic{figure}}
\renewcommand{\theequation}{S\arabic{equation}}

\begin{document}

\title{Supplementary Material for ``Repairing the Teacher for Knowledge
Distillation via Causally Screened Channel Edits''}

\author{Anonymous Authors}

\markboth{IEEE Transactions on Pattern Analysis and Machine Intelligence,~Vol.~XX, No.~XX, September~2026}%
{Anonymous Authors: Supplementary Material}

\maketitle

This supplement gives the experimental settings and further results
of the paper. Its sections, tables and equations carry the prefix
S; other numbers refer to the main paper.

\raggedbottom
\section{Datasets}
\label{app:datasets}

Table~\ref{tab:datasets} gives the sizes and splits of the thirteen
datasets, and Table~\ref{tab:dataset-classes} the classes of the ten
with at most eleven classes.
We obtained each dataset from its official distribution and checked
every downloaded archive against a fixed checksum.

\begin{table*}[!t]
\centering
\caption{Dataset characteristics and split sizes in our experiments.
Counts are numbers of images. The nominal imbalance factor $\rho$
controls class subsampling in the training and validation splits;
at $\rho=1$, no images are removed. Lastly, the test split is never resampled.}
\label{tab:datasets}
\footnotesize
\setlength{\tabcolsep}{4.5pt}
\begin{tabular}{@{}llrlrrrrrrl@{}}
\toprule
& & & & \multicolumn{3}{c}{$\rho=1$} & \multicolumn{3}{c}{$|\mathcal{D}_{\mathrm{tr}}|$ at} & \\
\cmidrule(lr){5-7}\cmidrule(lr){8-10}
Dataset & Domain & $C$ & Native size & $|\mathcal{D}_{\mathrm{tr}}|$ & $|\mathcal{D}_{\mathrm{val}}|$ & $|\mathcal{D}_{\mathrm{test}}|$ & $\rho=10$ & $\rho=50$ & $\rho=100$ & Validation split \\
\midrule
CIFAR-10      & Natural        & 10  & $32\times32$              & 45{,}000 & 5{,}000  & 10{,}000 & 18{,}391 & 12{,}601 & 11{,}167 & 10\% of train \\
CIFAR-100     & Natural        & 100 & $32\times32$              & 45{,}000 & 5{,}000  & 10{,}000 & 17{,}659 & 11{,}384 & 9{,}806  & 10\% of train \\
CINIC-10      & Natural        & 10  & $32\times32$              & 90{,}000 & 90{,}000 & 90{,}000 & 36{,}780 & 25{,}199 & 22{,}335 & Official \\
STL-10        & Natural        & 10  & $96\times96$              & 4{,}500  & 500      & 8{,}000  & 1{,}839  & 1{,}259  & 1{,}118  & 10\% of train \\
Caltech-101   & Natural        & 101 & Variable                  & 6{,}073  & 868      & 1{,}736  & 3{,}292  & 2{,}606  & 2{,}426  & 70/10/20 split \\
\midrule
SVHN          & House numbers  & 10  & $32\times32$              & 65{,}931 & 7{,}326  & 26{,}032 & 30{,}463 & 21{,}329 & 18{,}847 & 10\% of train \\
Fashion-MNIST & Clothing       & 10  & $28\times28$, grayscale   & 54{,}000 & 6{,}000  & 10{,}000 & 22{,}069 & 15{,}120 & 13{,}401 & 10\% of train \\
GTSRB         & Traffic signs  & 43  & Variable                  & 23{,}976 & 2{,}664  & 12{,}630 & 12{,}070 & 8{,}721  & 7{,}770  & 10\% of train \\
EuroSAT       & Satellite      & 10  & $64\times64$              & 21{,}600 & 2{,}700  & 2{,}700  & 9{,}118  & 6{,}408  & 5{,}732  & 80/10/10 split \\
\midrule
PathMNIST     & Colon pathology & 9  & $28\times28$              & 89{,}996 & 10{,}004 & 7{,}180  & 36{,}074 & 24{,}664 & 21{,}870 & Official \\
BloodMNIST    & Blood cells    & 8   & $28\times28$              & 11{,}959 & 1{,}712  & 3{,}421  & 4{,}647  & 3{,}099  & 2{,}714  & Official \\
DermaMNIST    & Dermatoscopy   & 7   & $28\times28$              & 7{,}007  & 1{,}003  & 2{,}005  & 1{,}722  & 874      & 707      & Official \\
OrganAMNIST   & Abdominal CT   & 11  & $28\times28$, grayscale   & 34{,}561 & 6{,}491  & 17{,}778 & 11{,}131 & 6{,}517  & 5{,}481  & Official \\
\bottomrule
\end{tabular}
\end{table*}

\begin{table}[t]
\caption{Classes of the ten datasets with at most eleven classes, in
label order. Bold classes are target classes at every $\rho>1$.}
\label{tab:dataset-classes}
\centering
\footnotesize
\setlength{\tabcolsep}{3pt}
\renewcommand{\arraystretch}{1.1}
\begin{tabular}{@{}>{\raggedright\arraybackslash}p{2cm}%
>{\raggedright\arraybackslash}p{0.73\columnwidth}@{}}
\toprule
Dataset & Classes, from label 0 \\
\midrule
CIFAR-10, CINIC-10 & airplane, automobile, bird, cat, deer,
\textbf{dog, frog, horse, ship, truck} \\
STL-10 & airplane, bird, car, cat, deer,
\textbf{dog, horse, monkey, ship, truck} \\
SVHN & 0, 1, 2, 3, 4, \textbf{5, 6, 7, 8, 9} \\
Fashion-MNIST & T-shirt/top, trouser, pullover, dress, coat,
\textbf{sandal, shirt, sneaker, bag, ankle boot} \\
EuroSAT & annual crop, forest, herbaceous vegetation, highway,
industrial, \textbf{pasture, permanent crop, residential, river,
sea/lake} \\
PathMNIST & adipose, background, debris, lymphocytes, \textbf{mucus,
smooth muscle, normal colon mucosa, cancer-associated stroma,
colorectal adenocarcinoma epithelium} \\
BloodMNIST & basophil, eosinophil, erythroblast, immature granulocyte,
\textbf{lymphocyte, monocyte, neutrophil, platelet} \\
DermaMNIST & actinic keratosis and intraepithelial carcinoma$^{a}$,
basal cell carcinoma, benign keratosis-like lesion,
\textbf{dermatofibroma, melanoma}, melanocytic nevus$^{b}$,
\textbf{vascular lesion} \\
OrganAMNIST & bladder, left femur, right femur$^{a}$, \textbf{heart},
left kidney, right kidney$^{b}$, liver, \textbf{left lung, right lung,
pancreas, spleen} \\
\bottomrule
\multicolumn{2}{@{}l}{$^{a}$Target only at $\rho=10$.\quad
$^{b}$Target only at $\rho=50$ and $\rho=100$.} \\
\end{tabular}
\end{table}

\textbf{Natural images.} CIFAR-10 and
CIFAR-100~\cite{krizhevsky2009learning} each contain 60{,}000 color
images of $32\times32$ pixels in 10 and 100 classes, with 50{,}000
training and 10{,}000 test images. CINIC-10~\cite{darlow2018cinic} combines
CIFAR-10 with downsampled ImageNet images of the same ten classes,
giving 270{,}000 images split equally into training, validation and
test sets. STL-10~\cite{coates2011analysis} has $96\times96$ images of
ten classes taken from ImageNet, with 500 labelled training and 800
test images per class; we do not use its 100{,}000 unlabelled images.
Caltech-101~\cite{feifei2004learning} contains photographs of objects in
101 categories at varying resolutions. Most categories hold about 50
images and the largest about 800, and
we exclude its background category.

\textbf{Other visual domains.} SVHN~\cite{netzer2011reading} contains
$32\times32$ crops of house-number digits from Google Street View
images, with 73{,}257 training and 26{,}032 test images; we do not use
its 531{,}131 additional training images.
Fashion-MNIST~\cite{xiao2017fashion} has 70{,}000 grayscale
$28\times28$ images of Zalando clothing articles in ten classes, split
into 60{,}000 training and 10{,}000 test images.
GTSRB~\cite{stallkamp2012man} contains photographs of German traffic
signs in 43 classes, from 15 to 250 pixels on a side; we use its
26{,}640-image training archive and 12{,}630-image test set.
EuroSAT~\cite{helber2019eurosat} has 27{,}000 Sentinel-2
patches of $64\times64$ pixels in ten land-use and land-cover classes,
of which we use the RGB version.

\textbf{Medical images.} The four medical datasets come from
MedMNIST~v2~\cite{yang2023medmnist}, which standardizes public medical
collections as $28\times28$ images with official training, validation
and test splits. PathMNIST holds patches of colorectal cancer histology
in nine tissue types, and BloodMNIST microscope images of individual
normal blood cells in eight types. DermaMNIST holds dermatoscopic images
of pigmented skin lesions in seven diagnostic categories, derived from
HAM10000, and OrganAMNIST axial slices of abdominal computed tomography
in eleven organ classes.

\textbf{Training split.} Each dataset has disjoint training, validation
and test sets, denoted $\mathcal{D}_{\mathrm{tr}}$,
$\mathcal{D}_{\mathrm{val}}$ and $\mathcal{D}_{\mathrm{test}}$.
Every teacher and student is trained on $\mathcal{D}_{\mathrm{tr}}$.
This set also provides the repair pool,
$\mathcal{D}_{\mathrm{rep}}=\mathcal{D}_{\mathrm{tr}}$, from which the
hard and preservation sets are drawn. For GTSRB~\cite{stallkamp2012man},
we use the 26{,}640-image competition training
archive.
For STL-10, we use the 5{,}000 labelled training images, and for SVHN we
do not use its extra training set.

\textbf{Validation split.} We use $\mathcal{D}_{\mathrm{val}}$ to select
checkpoints and repair candidates. CINIC-10 and the four MedMNIST
datasets provide official validation sets. For CIFAR-10, CIFAR-100,
STL-10, SVHN, Fashion-MNIST and GTSRB, we hold out a stratified 10\% of
the official training images. Caltech-101 and EuroSAT have no fixed
test partition, so we define stratified training, validation and test
splits for this study. Caltech-101 has 101 categories, most with about
50 images~\cite{feifei2004learning}; we assign 70/10/20 so that
its test set includes more images per class. EuroSAT has 27{,}000
images in 10 classes, and its original evaluation used 80\% for
training~\cite{helber2019eurosat}. We retain that training fraction
and divide the remaining 20\% equally between validation and test,
giving 80/10/10 and a 2{,}700-image test set. Within each dataset and
imbalance setting, every model uses the same resulting splits.

\textbf{Test split.} We keep $\mathcal{D}_{\mathrm{test}}$ sealed until
the teacher and every student checkpoint of a run have been fixed.
Official test sets are used unchanged, including the 12{,}630 GTSRB
and 8{,}000 STL-10 images. Caltech-101 and EuroSAT use
the partitions above. Test images are never resampled.

\textbf{Imbalance profiles.} Eq.~\eqref{M-eq:imbalance-profile} of the
main paper sets the imposed decay through $\rho$, but the ratio between
the largest and smallest class counts also depends on the native counts
and rounding. At $\rho=1$ no image is removed, so a dataset keeps its
native imbalance: DermaMNIST, for example, has between 80 and 4{,}693
training images per class. At $\rho=100$ the rarest training class holds
as few as one image (Caltech-101 and DermaMNIST).

\textbf{Preprocessing.} All images are used as $32\times32$ RGB tensors,
with grayscale images replicated across three channels and normalized per
channel with fixed constants: each dataset's own statistics for
CIFAR-100, CINIC-10 and Fashion-MNIST, the ImageNet statistics for
STL-10, Caltech-101, EuroSAT and GTSRB, and a mean and standard
deviation of 0.5 for the rest. For CIFAR-10, CIFAR-100,
CINIC-10 and SVHN, which are natively $32\times32$, and for Fashion-MNIST
and the MedMNIST datasets, which are resized from $28\times28$, training
images are randomly cropped to $32\times32$ after 4-pixel zero padding.
For the variable- and higher-resolution datasets (STL-10, Caltech-101,
EuroSAT and GTSRB), training uses a random resized crop to $32\times32$
covering 85--100\% of the image with aspect ratio 0.9--1.1, and
evaluation resizes the shorter side to 36 pixels and takes the central
$32\times32$ crop. Random horizontal flips are applied only to CIFAR-10,
CIFAR-100 and CINIC-10, where mirroring preserves the label; they are
omitted for digits, traffic signs, medical images and the remaining
datasets. Evaluation images are never augmented.

\textbf{Classes and target classes.}
Eq.~\eqref{M-eq:imbalance-profile} of the main paper thins classes in label order, so the
order decides which classes become rare.
Table~\ref{tab:dataset-classes} lists the classes of the ten datasets
with at most eleven classes in that order and marks the target classes
at $\rho>1$. When the native classes are of similar size, the target
classes are the later half of the order. Where native sizes differ, a
class that is small from the start can join the target set at some
imbalance factors only, as the footnotes of
Table~\ref{tab:dataset-classes} mark. For CIFAR-100, Caltech-101
and GTSRB we keep the standard label order, alphabetical for the first
two and by class number for GTSRB.
Their target sets hold 50, 51 and 22 classes. For CIFAR-100 these are
always the last 50 labels, whereas for Caltech-101 and GTSRB 45 to 47
and 20 of them lie in the later half of the order.

\section{Models and Training Settings}
\label{supp:training}

\textbf{Model architectures.} The design uses three teachers, CNN-T,
ResNet-18 and ResNet-34, and one student, CNN-S;
Table~\ref{tab:architectures} lists their structure, size and
compute. All four take $32\times32$ RGB inputs. CNN-T and CNN-S share
one design. Each has three stages, and each stage applies two
$3\times3$ convolutions, each followed by batch normalization and
ReLU, and ends with a $2\times2$ max pool that halves the spatial
resolution. The final feature map is average-pooled to $2\times2$ and
passed to a linear classifier. CNN-S halves CNN-T's width at every
stage. ResNet-18 and ResNet-34~\cite{he2016deep} stack basic residual
blocks, each with two $3\times3$ convolutions, batch normalization and
a shortcut connection, in four stages; the first block of each stage
after the first halves the spatial resolution. Both ResNets use a
$3\times3$, stride-one stem without initial max pooling, which keeps
the full $32\times32$ resolution in the first stage, and end with
global average pooling and a linear classifier. They share their
stage widths and differ only in depth. Stage widths in the table give
the output channels in order; stage depth means convolutions per stage
for the CNNs and residual blocks per stage for the ResNets. For each
teacher, the table also gives the number of channels that the repair
searches: the output channels of its last two stages.

\begin{table}[!t]
\caption{Teacher and student architectures. Parameters and FLOPs are
for a ten-class head and one $32\times32$ image, counting two FLOPs per
multiply--accumulate in convolution and linear layers.}
\label{tab:architectures}
\centering
\footnotesize
\setlength{\tabcolsep}{3.5pt}
\renewcommand{\arraystretch}{1.15}
\begin{tabular}{@{}lcccc@{}}
\toprule
& \multicolumn{3}{c}{Teachers} & Student \\
\cmidrule(lr){2-4}\cmidrule(l){5-5}
& CNN-T & ResNet-18 & ResNet-34 & CNN-S \\
\midrule
Stage widths & $(32,64,128)$ & \multicolumn{2}{c}{$(64,128,256,512)$}
& $(16,32,64)$ \\
Stage depth & $2$ & $(2,2,2,2)$ & $(3,4,6,3)$ & $2$ \\
Params (M) & 0.29 & 11.2 & 21.3 & 0.075 \\
FLOPs (M) & 77 & 1{,}111 & 2{,}319 & 20 \\
\multicolumn{5}{@{}l}{\textit{Ratio to CNN-S}} \\
\quad Params & $3.9\times$ & $149\times$ & $284\times$ & $1\times$ \\
\quad FLOPs & $3.9\times$ & $56\times$ & $117\times$ & $1\times$ \\
Searched channels & 192 & 768 & 768 & -- \\
\bottomrule
\end{tabular}
\end{table}

\textbf{Training settings.} Table~\ref{tab:training-parameters}
summarizes the optimization and distillation parameters. A separate
teacher is trained from scratch for each dataset--profile and
teacher-architecture setting, and CNN-S is trained with the same recipe
whichever teacher it learns from.

\begin{table*}[!t]
\caption{Training and distillation settings. Alternative KD losses use
the original teacher.}
\label{tab:training-parameters}
\centering
\footnotesize
\setlength{\tabcolsep}{4pt}
\renewcommand{\arraystretch}{1.08}
\begin{tabular}{@{}L{0.27\textwidth}L{0.67\textwidth}@{}}
\toprule
Parameter & Setting \\
\midrule
Teacher training & Cross-entropy from a fresh initialization \\
Optimizer & SGD; batch size 128; initial learning rate 0.05; momentum 0.9; weight decay $5\times10^{-4}$ \\
Epochs by imbalance factor & $(81,49,71,81)$ for $\rho=(1,10,50,100)$; cosine decay \\
Checkpoint selection & Highest validation accuracy for teacher and student training \\
Student initializations & Three per condition, paired across teacher conditions and KD losses \\
\midrule
Classical KD & Frozen original or repaired teacher; $\tau=4$; cross-entropy and response-KL weights 0.5 each; five-epoch response warm-up; no feature loss \\
LoCa & Cross-entropy and response weights 0.5 each; calibration coefficient 0.95 \\
DKD and RLD & Cross-entropy weight 1; response weight 1; internal component weights 1 and 8 \\
DKD and RLD, lower response & Cross-entropy weight 1; internal component weights 1 and 8; response weight set per teacher from training histories (see text), DKD/RLD: 0.15/0.20 with CNN-T, 0.298/0.516 with ResNet-18, and 0.290/0.539 with ResNet-34 \\
\bottomrule
\end{tabular}
\end{table*}

We ran all experiments on 16~GB NVIDIA GeForce RTX 5060 Ti
GPUs, using bfloat16 mixed precision with PyTorch~2.14 and
CUDA~13.2. All training used the same GPU model and precision, which
supported every teacher and the student at the batch sizes listed in
Table~\ref{tab:training-parameters}.

\section{Teacher Repair Settings}
\label{supp:repair}

\textbf{Teacher repair.} Table~\ref{tab:repair-parameters} gives the
values used for the parameters introduced in Section~\ref{M-sec:repair} of the main paper.
Localization examines the last two convolutional stages of the
teacher. For sparse classes, the companion and preservation pools use
the fallback rules described in Section~\ref{M-sec:repair}; the table
reports the requested companion count and preservation cap. Candidates
are selected on the validation set as in Step 6, with a selection
guardrail of $\eta_{\mathrm{sel}}=0.005$. If no candidate passes, the
original teacher is retained as the no-change outcome, and the
original-teacher student is reused rather than retrained.

\begin{table}[!t]
\caption{Parameters used for teacher localization and repair.}
\label{tab:repair-parameters}
\centering
\footnotesize
\setlength{\tabcolsep}{3pt}
\renewcommand{\arraystretch}{1.08}
\begin{tabular}{@{}L{0.39\columnwidth}L{0.53\columnwidth}@{}}
\toprule
Component & Setting \\
\midrule
Hard set and companions & $\xi=0.20$; $n_{\mathrm{comp}}=5$ \\
Stable channel screen & $B_{\mathrm{boot}}=20$; $\kappa=0.25$; $\pi_{\min}=0.70$ \\
Causal preservation set & Equal count per usable class, capped at 64 \\
Causal guardrail & $\eta_{\mathrm{acc}}=0.005$ \\
Selection guardrail & $\eta_{\mathrm{sel}}=0.005$ \\
Candidate grid & $n_{\mathrm{edit}}\in\{2,4,8\}$ for every teacher; $\{8,16,32\}$ for the scaled ResNet grid; learning rates 0.001 and 0.005 \\
Repair optimization & 20 epochs; 640 draws per stream per epoch; batch size 64 per stream; SGD momentum 0.9, zero weight decay and cosine schedule \\
Preservation objective & $\tau_{\mathrm r}=4$; $\lambda_{\mathrm P}=1$; $\lambda_{\mathrm F}=0.25$ \\
\bottomrule
\end{tabular}
\end{table}

\section{Evaluation Metrics}
\label{supp:metrics}

This section defines the metrics of Section~\ref{M-sec:evaluation-metrics}
of the main paper on the test set
$\mathcal{D}_{\mathrm{test}}=\{(x_i,y_i)\}_{i=1}^{N}$, where $x_i$ is
an image and $y_i$ its class label. For a model $M$, let
$\widehat y_i(M)=\argmax_c p_{M,c}^{(1)}(x_i)$ be its predicted label.
The $(a,b)$ entry of its confusion matrix is
\begin{equation}
\label{eq:confusion-entry}
Q_{ab}(M)=\sum_{i=1}^{N}
\mathbf{1}[y_i=a]\cdot\mathbf{1}[\widehat y_i(M)=b].
\end{equation}
Thus, $Q_{cc}(M)$ counts the class-$c$ images predicted correctly,
whereas an off-diagonal entry $Q_{ab}(M)$ counts class-$a$ images
mistaken for class $b$. We compute accuracy, precision
and recall from these counts. Overall accuracy is
\begin{equation}
\label{eq:overall-accuracy}
\operatorname{Acc}(M)=\frac{1}{N}\sum_{c=0}^{C-1}Q_{cc}(M).
\end{equation}
For class $c$, precision and recall are respectively
\begin{equation}
\label{eq:class-precision}
P_c(M)=\frac{Q_{cc}(M)}{\sum_a Q_{ac}(M)},
\end{equation}
and
\begin{equation}
\label{eq:class-recall}
R_c(M)=\frac{Q_{cc}(M)}{\sum_b Q_{cb}(M)}.
\end{equation}
If a class receives no predictions, its precision is set to zero.
We average class F1 scores to obtain
\begin{equation}
\label{eq:macro-f1}
F_{1,\mathrm{macro}}(M)=\frac{1}{C}
\sum_{c=0}^{C-1}\frac{2P_c(M)R_c(M)}{P_c(M)+R_c(M)},
\end{equation}
where a class with zero precision and recall contributes zero.

The primary class-sensitive measure, target-class recall
$R_{\mathrm{tar}}$, averages $R_c$ over $\mathcal C_{\mathrm{tar}}$
(Eq.~\eqref{M-eq:target-recall} of the main paper).

We also examine the predicted probabilities. Let
$p_{i,c}=p_{M,c}^{(1)}(x_i)$ denote the probability that model $M$
assigns to class $c$ for image $x_i$, and let
$N_c=\sum_{i=1}^{N}\mathbf{1}[y_i=c]$ be the number of test images
whose true label is $c$. Overall negative log-likelihood (NLL), Brier
score, and target-balanced NLL are defined separately. Overall NLL is
\begin{equation}
\label{eq:overall-nll}
\operatorname{NLL}(M)
=-\frac{1}{N}\sum_{i=1}^{N}\log p_{i,y_i}.
\end{equation}
The multiclass Brier score is
\begin{equation}
\label{eq:brier-score}
\operatorname{Brier}(M)
=\frac{1}{N}\sum_{i=1}^{N}\sum_{c=0}^{C-1}
\bigl(p_{i,c}-\mathbf{1}[y_i=c]\bigr)^2.
\end{equation}
Finally, target-balanced NLL is
\begin{equation}
\label{eq:target-balanced-nll}
\operatorname{NLL}_{\mathrm{tar}}(M)
=-\frac{1}{|\mathcal C_{\mathrm{tar}}|}
\sum_{c\in\mathcal C_{\mathrm{tar}}}
\frac{1}{N_c}\sum_{i:y_i=c}\log p_{i,c}.
\end{equation}
The last measure weights target classes equally and breaks ties in
validation-based repair selection. Smaller values are better for all
three measures. To assess confidence calibration, we place each
prediction's maximum probability in one of 15 equal-width bins. If
$n_b$, $a_b$ and $q_b$ are the count, accuracy and mean confidence of
bin $b$, respectively, then
\begin{equation}
\label{eq:ece-metric}
\operatorname{ECE}_{15}(M)=\sum_{b=1}^{15}
\frac{n_b}{N}|a_b-q_b|,
\end{equation}
with empty bins contributing zero. We also assess whether confidence
identifies correct predictions. Order the $N$
images by decreasing maximum probability, write their correctness
indicators as $u_{(1)},\ldots,u_{(N)}$, and compute the area under the
risk--coverage curve (AURC) as
\begin{equation}
\label{eq:aurc-metric}
\operatorname{AURC}(M)=\frac{1}{N}\sum_{k=1}^{N}
\left(1-\frac{1}{k}\sum_{j=1}^{k}u_{(j)}\right).
\end{equation}
Tied confidence values are averaged over their possible orders.
Selective accuracy at 80\% or 90\% coverage is the accuracy among the
most confident $\lceil 0.8N\rceil$ or $\lceil 0.9N\rceil$ predictions.
Lower ECE and AURC indicate better calibration and selective
performance, respectively.

\section{Alternative-Loss Settings}
\label{app:loss-settings}

Here we state the alternative losses as implemented and how the lower
weights of DKD and RLD were set. Every loss, classical KD included,
trains the student with
\begin{equation}
\label{eq:alternative-student-loss}
\mathcal L=\lambda_{\mathrm{CE}}\,
\mathcal L_{\mathrm{CE}}\!\left(y,p_S^{(1)}(x)\right)
+w_e\lambda_{\mathrm R}\,\mathcal R(x,y),
\end{equation}
where $\mathcal R$ is the method's response term and $w_e$ is the
warm-up factor of Eq.~\eqref{M-eq:experimental-kd-loss} in the main paper. All four losses
use $\tau=4$ and no feature term. Classical KD and LoCa set
$\lambda_{\mathrm R}=0.5$ and $\lambda_{\mathrm{CE}}=1-\lambda_{\mathrm R}$,
as in the main paper, whereas DKD and RLD keep
$\lambda_{\mathrm{CE}}=1$ and use either the published
$\lambda_{\mathrm R}=1$ or the lower weight described below.

\begin{table}[t]
\caption{Response weights of DKD and RLD, and the median ratio of the
weighted response term to the weighted cross-entropy over all 52
settings.}
\label{tab:loss-weights}
\centering
\footnotesize
\setlength{\tabcolsep}{4pt}
\renewcommand{\arraystretch}{1.08}
\begin{tabular}{@{}lccc@{}}
\toprule
& CNN-T & ResNet-18 & ResNet-34 \\
\midrule
\multicolumn{4}{@{}l}{\textit{Lower response weight $\lambda_{\mathrm R}$}} \\
DKD & 0.15 & 0.298 & 0.290 \\
RLD & 0.20 & 0.516 & 0.539 \\
\midrule
\multicolumn{4}{@{}l}{\textit{Median ratio of response term to cross-entropy}} \\
Classical KD & 4.71 & 8.06 & 7.07 \\
LoCa & 4.82 & 8.03 & 7.10 \\
DKD, published weight & 27.85 & 27.02 & 24.35 \\
DKD, lower weight & 5.25 & 9.72 & 8.20 \\
RLD, published weight & 22.07 & 15.61 & 13.12 \\
RLD, lower weight & 5.20 & 9.00 & 7.52 \\
\bottomrule
\end{tabular}
\end{table}

\textbf{Response terms.} The response term of classical KD is
$\mathcal R=\tau^2D_{\mathrm{KL}}(p_T^{(\tau)}\|p_S^{(\tau)})$, and
LoCa~\cite{yang2024loca} uses the same term with a calibrated teacher
distribution. On an image where the teacher ranks another class above
the ground truth $y$, every non-target probability is multiplied by
$0.95/(1-p_{T,y}^{(\tau)}+\max_c p_{T,c}^{(\tau)})$ and $y$ receives
the remaining mass, which makes it the most probable class; on the
other images the teacher's distribution is unchanged.
DKD~\cite{zhao2022decoupled} uses $\mathcal R=\tau^2\mathcal
L_{\mathrm{DKD}}$ from Eq.~\eqref{M-eq:dkd-loss} of the main paper, with
$\lambda_{\mathrm{TC}}=1$ and $\lambda_{\mathrm{NC}}=8$.
RLD~\cite{sun2025knowledge} also has two terms, weighted 1 and 8. The
first matches the student's binary distribution for $y$ to the
teacher's binary distribution for its own top class, at temperature 1.
The second applies $\tau^2D_{\mathrm{KL}}$ at temperature $\tau$ over
only the classes that the teacher ranks below $y$. On an image that the
teacher classifies correctly, the two terms therefore reduce to DKD's
target and non-target terms, apart from the temperature of the first.

\textbf{Lower response weight.} The lower weight rescales
$\lambda_{\mathrm R}$ so that the response term relates to the
cross-entropy as it does in classical KD. For each run, we take the
ratio $\lambda_{\mathrm R}\mathcal R/(\lambda_{\mathrm{CE}}\mathcal
L_{\mathrm{CE}})$ of the two terms' epoch means, find its median over
the epochs after the warm-up, and then take the median over runs. The
lower weight is classical KD's median divided by the method's median at
$\lambda_{\mathrm R}=1$, so it depends on training histories alone. For
the ResNet teachers we computed it over all 52 settings. For CNN-T we
fixed it earlier, on its 23 repaired settings, where the medians were
4.42 for classical KD, 29.02 for DKD and 21.59 for RLD. This gave 0.152
and 0.205, which we rounded to 0.15 and 0.20.
Table~\ref{tab:loss-weights} gives the weights and the medians over all
52 settings. At the lower weights, the DKD and RLD ratios came to 1.06
to 1.21 times that of classical KD, against 1.9 to 5.9 times at the
published weights. LoCa kept the balance of classical KD, since it
changes only the teacher's targets.

\section{Dataset-Level Student Performance\\by Teacher}
\label{app:dataset-student-outcomes}

Table~\ref{tab:dataset-student-outcomes} counts how often each repaired
teacher improved, left unchanged, or decreased its CNN-S student's
test performance relative to the matched original-teacher condition,
under both channel budgets.
Each dataset contributes one result per imbalance factor. Counts use the sign of the student difference averaged over
three matched initializations. A zero difference includes settings in
which no teacher repair was selected and the original student was
reused.

\begin{table}[t]
\caption{CNN-S outcomes by dataset for each teacher and channel budget,
counted as improved/unchanged/decreased settings over the four
imbalance factors.}
\label{tab:dataset-student-outcomes}
\centering
\footnotesize
\renewcommand{\arraystretch}{1.05}
\setlength{\tabcolsep}{3pt}
\begin{tabular}{@{}lccccc@{}}
\toprule
& CNN-T & \multicolumn{2}{c}{ResNet-18} & \multicolumn{2}{c}{ResNet-34} \\
\cmidrule(lr){2-2}\cmidrule(lr){3-4}\cmidrule(l){5-6}
Dataset & 2/4/8 & 2/4/8 & 8/16/32 & 2/4/8 & 8/16/32 \\
\midrule
\multicolumn{6}{@{}l}{\textit{Test accuracy}} \\
CIFAR-10 & 2/1/1 & 1/1/2 & 2/1/1 & 3/1/0 & 3/1/0 \\
CIFAR-100 & 0/4/0 & 2/1/1 & 2/2/0 & 2/1/1 & 0/3/1 \\
CINIC-10 & 2/2/0 & 3/0/1 & 2/1/1 & 3/0/1 & 4/0/0 \\
STL-10 & 1/3/0 & 0/4/0 & 0/4/0 & 1/3/0 & 1/3/0 \\
Caltech-101 & 1/3/0 & 0/3/1 & 0/3/1 & 0/4/0 & 0/4/0 \\
SVHN & 2/2/0 & 0/3/1 & 0/3/1 & 0/3/1 & 0/3/1 \\
Fashion-MNIST & 3/1/0 & 1/3/0 & 1/3/0 & 3/1/0 & 2/2/0 \\
GTSRB & 0/4/0 & 0/4/0 & 1/3/0 & 0/4/0 & 0/3/1 \\
EuroSAT & 1/2/1 & 1/3/0 & 1/3/0 & 2/2/0 & 2/2/0 \\
PathMNIST & 2/1/1 & 0/4/0 & 0/4/0 & 0/2/2 & 1/2/1 \\
BloodMNIST & 3/1/0 & 1/3/0 & 1/3/0 & 0/3/1 & 0/4/0 \\
DermaMNIST & 1/2/1 & 2/1/1 & 1/1/2 & 0/4/0 & 0/4/0 \\
OrganAMNIST & 1/3/0 & 0/4/0 & 0/4/0 & 1/3/0 & 0/3/1 \\
\midrule
All 52 settings & 19/29/4 & 11/34/7 & 11/35/6 & 15/31/6 & 13/34/5 \\
\midrule
\multicolumn{6}{@{}l}{\textit{Target-class macro recall}} \\
CIFAR-10 & 2/1/1 & 0/1/3 & 1/1/2 & 3/1/0 & 3/1/0 \\
CIFAR-100 & 0/4/0 & 3/1/0 & 2/2/0 & 2/1/1 & 0/3/1 \\
CINIC-10 & 2/2/0 & 3/0/1 & 2/1/1 & 2/0/2 & 3/0/1 \\
STL-10 & 1/3/0 & 0/4/0 & 0/4/0 & 1/3/0 & 1/3/0 \\
Caltech-101 & 1/3/0 & 0/3/1 & 1/3/0 & 0/4/0 & 0/4/0 \\
SVHN & 2/2/0 & 0/3/1 & 0/3/1 & 0/3/1 & 0/3/1 \\
Fashion-MNIST & 3/1/0 & 1/3/0 & 1/3/0 & 3/1/0 & 2/2/0 \\
GTSRB & 0/4/0 & 0/4/0 & 1/3/0 & 0/4/0 & 0/3/1 \\
EuroSAT & 2/2/0 & 1/3/0 & 1/3/0 & 2/2/0 & 2/2/0 \\
PathMNIST & 2/1/1 & 0/4/0 & 0/4/0 & 0/2/2 & 1/2/1 \\
BloodMNIST & 3/1/0 & 1/3/0 & 1/3/0 & 0/3/1 & 0/4/0 \\
DermaMNIST & 1/2/1 & 3/1/0 & 2/1/1 & 0/4/0 & 0/4/0 \\
OrganAMNIST & 1/3/0 & 0/4/0 & 0/4/0 & 1/3/0 & 1/3/0 \\
\midrule
All 52 settings & 20/29/3 & 12/34/6 & 12/35/5 & 14/31/7 & 13/34/5 \\
\bottomrule
\end{tabular}
\end{table}

Under the fixed budget, the accuracy counts are 19/29/4 for CNN-T,
11/34/7 for ResNet-18 and 15/31/6 for ResNet-34, and the target-recall
counts are 20/29/3, 12/34/6 and 14/31/7. Under the scaled budget, the
two ResNet teachers reach 11/35/6 and 13/34/5 in accuracy and 12/35/5
and 13/34/5 in target recall, close to their fixed-budget counts.

\section{Additional Metrics of the\\Selected Repairs}
\label{supp:secondary-results}

Table~\ref{tab:secondary-results} gives the mean changes in the
additional metrics of Section~\ref{supp:metrics} over each teacher's
selected repairs. Student differences are averaged first over the three
matched initializations and then across settings. Lower NLL, Brier
score, ECE and AURC indicate better performance, as do higher F1 and
selective accuracy. For every teacher, the mean changes favored both
the repaired teacher and its students in all eight metrics. The
improvements were smaller for the ResNet teachers, in line with their
lower recall gains, although individual settings can differ from the
mean.

\begin{table}[!t]
\caption{Mean repaired-minus-original changes in the additional
metrics over each teacher's selected repairs (23, 18 and 21
settings), with CNN-S as the student. F1, ECE and selective accuracy
are in pp, and the other rows use their native scales.}
\label{tab:secondary-results}
\centering
\footnotesize
\setlength{\tabcolsep}{4pt}
\renewcommand{\arraystretch}{1.1}
\begin{tabular}{@{}lrrr@{}}
\toprule
Metric & CNN-T & ResNet-18 & ResNet-34 \\
\midrule
\multicolumn{4}{@{}l}{\textit{Teacher}} \\
Macro F1 (pp) & $+1.06$ & $+0.62$ & $+0.29$ \\
Target-balanced NLL & $-0.155$ & $-0.119$ & $-0.055$ \\
Overall NLL & $-0.057$ & $-0.026$ & $-0.023$ \\
Brier score & $-0.016$ & $-0.007$ & $-0.005$ \\
ECE$_{15}$ (pp) & $-1.07$ & $-0.56$ & $-0.28$ \\
AURC & $-0.0055$ & $-0.0027$ & $-0.0016$ \\
Selective accuracy, 80\% (pp) & $+0.87$ & $+0.31$ & $+0.24$ \\
Selective accuracy, 90\% (pp) & $+0.91$ & $+0.36$ & $+0.29$ \\
\midrule
\multicolumn{4}{@{}l}{\textit{Student}} \\
Macro F1 (pp) & $+0.73$ & $+0.19$ & $+0.13$ \\
Target-balanced NLL & $-0.115$ & $-0.065$ & $-0.034$ \\
Overall NLL & $-0.040$ & $-0.005$ & $-0.009$ \\
Brier score & $-0.011$ & $-0.003$ & $-0.001$ \\
ECE$_{15}$ (pp) & $-0.77$ & $-0.26$ & $-0.13$ \\
AURC & $-0.0032$ & $-0.0012$ & $-0.0003$ \\
Selective accuracy, 80\% (pp) & $+0.50$ & $+0.17$ & $+0.05$ \\
Selective accuracy, 90\% (pp) & $+0.51$ & $+0.20$ & $+0.08$ \\
\bottomrule
\end{tabular}
\end{table}

\section{Notation}
\label{app:notation}

Table~\ref{tab:notation} lists the main symbols by where they are
introduced.

% A breakable table, so that it starts in the space after this section's
% text and continues in the next column under a repeated head.
\vspace{\intextsep}
\flowtablecaption{Main symbols.\label{tab:notation}}
{\footnotesize
\renewcommand{\arraystretch}{1.08}
\setlength{\tabcolsep}{3pt}
\tablefirsthead{\toprule Symbol & Meaning \\ \midrule}
\tablehead{\toprule Symbol & Meaning \\ \midrule}
\tabletail{\bottomrule}
\tablelasttail{\bottomrule}
\centering
\begin{supertabular}{@{}>{\raggedright\arraybackslash}p{0.30\columnwidth}%
>{\raggedright\arraybackslash}p{0.66\columnwidth}@{}}
% supertabular's row-height estimate is conservative; let the first part use
% the rest of its column.
\shrinkheight{-36pt}
\multicolumn{2}{@{}l}{\textit{Distillation (Section~\ref{M-sec:introduction})}} \\
$T$, $S$ & Teacher and student \\
$z_M(x)$, $p_M^{(\tau)}(x)$ & Logits of model $M$ and its softmax at temperature $\tau$ \\
$\mathcal Y$, $C$ & Set of class labels and its size \\
$\mathcal L_{\mathrm{KD}}$, $\alpha$ & KD objective and its cross-entropy weight \\
$\mathcal E(f)$, $\mathcal E_k(f)$ & Error rate of classifier $f$, overall and on class $k$ \\
$\Phi_\tau(S,T)$ & Fidelity gap, the expected $D_{\mathrm{KL}}(p_T^{(\tau)}\|p_S^{(\tau)})$ \\
$\gamma$, $\gamma_S$ & Teacher and student capacity \\
$\mathcal Y_U$ & Undistillable classes, on which KD does worse than cross-entropy alone \\
\midrule
\multicolumn{2}{@{}l}{\textit{Repair (Section~\ref{M-sec:repair})}} \\
$T_0=f_{\theta_0}$, $T_\theta$ & Original teacher and repair candidate, with $d_\theta$ parameters \\
$\mathcal D_{\mathrm{rep}}$ & Repair pool of $N$ labelled images, the training set \\
$\mathcal C_{\mathrm{tar}}$ & Target classes \\
$e_i$ & Error margin: strongest competitor's logit minus the true class's \\
$\mathcal H$, $H_c$, $\xi$ & Hard set, its class-$c$ indices and the hard-set fraction \\
$\mathcal P$, $\mathcal P_{\mathrm{causal}}$ & Preservation pool and its class-balanced subset for causal screening \\
$\mathcal N_i$, $n_{\mathrm{comp}}$ & Companion list of hard image $i$ and its length \\
$\mathcal L_{\mathrm{loc}}$, $\mathcal A_{\mathrm{loc}}$, $a=(\ell,k)$ & Searched stages, their channels and one channel \\
$A_\theta^{(\ell)}(x)$, $K_\ell$ & Feature map of stage $\ell$ and its number of channels \\
$G_{i,\ell,k}$, $D_{i,\ell,k}$, $\overline s_a$ & Gradient sensitivity, activation difference and consensus score of a channel \\
$B_{\mathrm{boot}}$, $\kappa$, $\pi_{\min}$ & Bootstrap resamples, retained fraction and minimum selection frequency \\
$\mathcal A_{\mathrm{stable}}$ & Channels that pass the stability screen \\
$T_0^{-a}$ & Original teacher with channel $a$ set to zero \\
$m(T;x,y)$ & True-class margin, equal to $-e_i$ for $T_0$ \\
$\eta_{\mathrm{acc}}$ & Causal guardrail, the largest accuracy drop allowed on $\mathcal P_{\mathrm{causal}}$ \\
$n_{\mathrm{edit}}$ & Channel budget, the number of channels edited \\
$\mathcal A_{\mathrm{edit}}$, $\Omega_{\mathrm{edit}}$ & Edited channels and their parameter coordinates \\
$\mathcal L_{\mathrm{rep}}$, $d_{\mathrm F}$ & Repair objective and its feature distance \\
$\tau_{\mathrm r}$, $\lambda_{\mathrm P}$, $\lambda_{\mathrm F}$ & Repair temperature and the weights of output and feature preservation \\
$\eta_{\mathrm{sel}}$ & Selection guardrail, the largest drop in validation accuracy allowed for a selected candidate \\
\midrule
\multicolumn{2}{@{}l}{\textit{Experiments (Sections~\ref{M-sec:experimental-design} and~\ref{M-sec:results})}} \\
$\mathcal D_{\mathrm{tr}}$, $\mathcal D_{\mathrm{val}}$, $\mathcal D_{\mathrm{test}}$ & Training, validation and test sets \\
$\rho$ & Nominal imbalance factor \\
$n_c$, $\widetilde n_c$ & Images kept in class $c$ and images available to it \\
$w_e$ & Warm-up factor of the response term in epoch $e$ \\
$\lambda_{\mathrm{CE}}$, $\lambda_{\mathrm R}$, $\mathcal R$ & Cross-entropy weight, response weight and response term (Section~\ref{app:loss-settings}) \\
$R_c$, $R_{\mathrm{tar}}$ & Recall of class $c$ and mean recall over $\mathcal C_{\mathrm{tar}}$ \\
$\Delta\mathrm{Acc}$, $\Delta R_{\mathrm{tar}}$ & Change when the teacher is repaired, in pp \\
\end{supertabular}\par
}

\FloatBarrier

\bibliographystyle{IEEEtran}
\bibliography{ref}

\end{document}
